\setcounter{section}{2}

  \zwischenueberschrift{Rotationsflächen}

Wir besprechen eine weitere Klasse von Flächen im Raum, die Rotationsflächen.

\bild{ \begin{center}
\includegraphics[width=5.5cm]{ \bildeinlesungsvg {Integral apl rot obsah1} {svg} }
\end{center}
\bildtext {} }

\bildlizenz { Integral apl rot obsah1.svg } {} {Pajs} {Tschechische Wikipedia} {gemeinfrei} {}

Es sei eine
\definitionsverweis {differenzierbare Kurve}{}{}
\maabbeledisp {\gamma} {]a,b[} {\R^2
} { t } {(x(t),y(t))
} {,}
mit

\mavergleichskette
{\vergleichskette
{ y(t)
}
{ \geq }{ 0
}
{  }{
}
{    }{
}
{   }{
}
}
{}{}{}
gegeben. Wir interessieren uns für die zugehörige
\definitionsverweis {Rotationsfläche}{}{,}
also die Teilmenge

\mathdisp {\left\{  (x(t), y(t) \cos \alpha , y(t) \sin \alpha)   \mid  t \in [a,b] , \, \alpha \in [0,2\pi]       \right\}} {  }

des
\mathl{\R^3}{,} die entsteht, wenn man die Trajektorie der Kurve um die $x$-Achse dreht. Wir setzen zusätzlich voraus, dass $\gamma$ einen Homöomorphismus auf sein Bild

\mavergleichskette
{\vergleichskette
{ M
}
{ = }{ \gamma \left(  ]a,b[  \right)
}
{  }{
}
{    }{
}
{   }{
}
}
{}{}{}
bewirkt. Insbesondere betrachten wir den Fall, wo
\maabb {f} {]a,b[} { \R_+
} {}
eine differenzierbare Funktion ist und es um die Rotationsfläche des Graphen geht.

__NOEDITSECTION__

\inputfaktbeweis
{Rotationsfläche/Differenzierbare positive Kurve/Graph/Tangentialraum/Fakt}
{Lemma}
{}
{

\faktsituation {Es sei $I$ ein
\definitionsverweis {offenes Intervall}{}{}
und
\maabb {f} { I } { \R_+
} {}
eine
\definitionsverweis {differenzierbare Funktion}{}{.}}
\faktuebergang {Dann besitzt die
\definitionsverweis {Rotationsfläche}{}{}
$Y$
\zusatzklammer {um die $x$-Achse} {} {}
zum
\definitionsverweis {Graphen}{}{}
von $f$ folgende Eigenschaften.}
\faktfolgerung {\aufzaehlungdrei{Es ist $Y$ die Nullstellenmenge zu
\maabbeledisp {h} {I \times \R \times \R} {\R
} { \left(  x   , \,   y  , \,  z                 \right) } { y^2+z^2-f(x)^2
} {.}
}{Die beschreibende Funktion $h$ ist in jedem Punkt von $Y$
\definitionsverweis {regulär}{}{.}
}{Der
\definitionsverweis {Tangentialraum}{}{}
von $Y$ in einem Punkt
\mathl{\left(  x   , \,   y  , \,  z                 \right)}{} ist
\zusatzklammer {bei

\mavergleichskettek
{\vergleichskettek
{ z
}
{ \neq }{ 0
}
{  }{
}
{    }{
}
{   }{
}
}
{}{}{}} {} {}

\mavergleichskettedisp
{\vergleichskette
{ T_PY
}
{ =} { \R \begin{pmatrix}  0  \\ z \\  -y   \end{pmatrix} + \R\begin{pmatrix}  z  \\ 0 \\  f(x)f'(x)   \end{pmatrix}
}
{ } {
}
{ } {
}
{ } {
}
}
{}{}{.}
}}
\faktzusatz {}
\faktzusatz {}

}
{

(1) ist klar. Die
\definitionsverweis {Jacobimatrix}{}{}
von $h$ ist
\mathl{2 \left(  -f(x)f'(x)  , \,   y  , \,  z                 \right)}{.} Wegen der Positivität von $f$ kann

\mavergleichskette
{\vergleichskette
{ y
}
{ = }{ z
}
{ = }{ 0
}
{    }{
}
{   }{
}
}
{}{}{}
nicht sein, daher ist $h$ auf $Y$ regulär. Die angegebenen Vektoren sind bei

\mavergleichskette
{\vergleichskette
{ z
}
{ \neq }{ 0
}
{  }{
}
{    }{
}
{   }{
}
}
{}{}{}
\definitionsverweis {linear unabhängig}{}{}
und gehören zum
\definitionsverweis {Kern}{}{}
der Jacobimatrix.

}

__NOEDITSECTION__

\inputfaktbeweis
{Rotationsfläche/Differenzierbare positive Kurve/Graph/Geodätische/Fakt}
{Lemma}
{}
{

\faktsituation {Es sei $I$ ein
\definitionsverweis {offenes Intervall}{}{}
und
\maabb {f} { I } { \R_+
} {}
eine
\definitionsverweis {differenzierbare Funktion}{}{}
und es sei $Y$ die
\definitionsverweis {Rotationsfläche}{}{}
\zusatzklammer {um die $x$-Achse} {} {}
zum
\definitionsverweis {Graphen}{}{}
von $f$. Es sei
\maabbeledisp {\gamma} { J } { I
} { t } { \gamma(t)
} {,}
eine bijektive differenzierbare Funktion derart, dass
\mathl{\left(  \gamma(t)  , \,  f(\gamma(t))                   \right)}{}
ein Parametrisierung des Graphen zu $f$ mit konstanter Norm der Geschwindigkeit ist.}
\faktuebergang {Dann gelten folgende Eigenschaften.}
\faktfolgerung {\aufzaehlungzwei {Zu jedem Punkt

\mavergleichskette
{\vergleichskette
{ \left(  x   , \,   y  , \,  z                 \right)
}
{ = }{ \left(  x   , \,   f(x) \cos \alpha  , \,   f(x) \sin \alpha                 \right)
}
{ \in }{ Y
}
{    }{
}
{   }{
}
}
{}{}{}
ist
\maabbeledisp {\theta} { J } { Y
} { t } { \left(  \gamma(t)   , \,  f( \gamma(t) ) \cos \alpha  , \,   f( \gamma(t) ) \sin \alpha                 \right)
} {,}
eine
\definitionsverweis {Geodätische}{}{.}
} {Eine Kreiskurve
\maabbeledisp {\theta} {\R} {I \times \R^2
} {\alpha} { \left(  x   , \,  f( x ) \cos \alpha  , \,   f(x ) \sin \alpha                 \right)
} {,}
ist genau dann eine Geodätische, wenn

\mavergleichskette
{\vergleichskette
{ f'(x)
}
{ = }{ 0
}
{  }{
}
{    }{
}
{   }{
}
}
{}{}{}
ist.
}}
\faktzusatz {}
\faktzusatz {}

}
{

\aufzaehlungzwei {Die Kurve $\theta$ bewegt sich in der Ebene
\mathl{\R  \begin{pmatrix}  1  \\ 0 \\  0   \end{pmatrix} + \R  \begin{pmatrix}  0  \\ \cos \alpha \\  \sin \alpha   \end{pmatrix}}{.} Nach einer Drehung können wir ohne Einschränkung

\mavergleichskette
{\vergleichskette
{ \alpha
}
{ = }{ 0
}
{  }{
}
{    }{
}
{   }{
}
}
{}{}{}
annehmen, die Kurve durchläuft also direkt den Graphen von $f$. Wegen der Bogenparametrisierung steht

\mavergleichskettedisp
{\vergleichskette
{ \theta'(t)
}
{ =} { \begin{pmatrix}  \gamma'(t) \\f'(\gamma(t)) \gamma'(t) \\  0   \end{pmatrix}
}
{ } {
}
{ } {
}
{ } {
}
}
{}{}{}
nach
[[Differenzierbare Kurve/Bogenparametrisiert/Zweite Ableitung/Senkrecht/Aufgabe|Aufgabe 2.3]]
senkrecht auf

\mavergleichskettedisp
{\vergleichskette
{ \theta^{\prime \prime} (t)
}
{ =} { \begin{pmatrix}  \gamma^{\prime \prime} (t) \\f'(\gamma(t)) \gamma^{\prime \prime} (t) +f^{\prime \prime}(\gamma(t))  \left(  \gamma^{\prime } (t)  \right)^2 \\  0   \end{pmatrix}
}
{ } {
}
{ } {
}
{ } {
}
}
{}{}{.}
Dies bedeutet, dass die Beschleunigung senkrecht auf dem Tangentialraum
\zusatzklammer {der neben der Kurvenableitung von $\begin{pmatrix}  0  \\ 0\\  1   \end{pmatrix}$ erzeugt wird} {} {}
steht und daher die tangentiale Beschleunigung gleich $0$ ist. Also liegt eine Geodätische vor.
} {Die Kreisbewegung spielt sich auf der durch das fixierte $x$ gegebenen Ebene ab, die Ableitung
\zusatzklammer {nach $\alpha$} {} {}
der gegebenen Kurve ist  $f(x) \begin{pmatrix}  0  \\ - \sin \alpha \\  \cos \alpha   \end{pmatrix}$  und die Beschleunigung der gegebenen Kurve ist gleich  $- f(x) \begin{pmatrix}  0  \\ \cos \alpha \\  \sin \alpha   \end{pmatrix}$. Die Beschleunigung ist damit innerhalb der Ebene senkrecht zur Geschwindigkeit der Kurve, die einen Tangentialvektor der Rotationsfläche bildet. Ein dazu linear unabhängiger Tangentialvektor ist durch $\begin{pmatrix}  1  \\f'(x)\\  0   \end{pmatrix}$ gegeben. Dieser steht aber nur bei

\mavergleichskette
{\vergleichskette
{ f'(x)
}
{ = }{ 0
}
{  }{
}
{    }{
}
{   }{
}
}
{}{}{}
stets senkrecht auf der Beschleunigung.
}

}

  \zwischenueberschrift{Das Normalenfeld}

\inputdefinition
{}
{

Es sei

\mavergleichskette
{\vergleichskette
{ W
}
{ \subseteq }{ \R^n
}
{  }{
}
{    }{
}
{   }{
}
}
{}{}{}
\definitionsverweis {offen}{}{,}
\maabb {h} { W } { \R
} {}
eine
\definitionsverweis {stetig differenzierbare Funktion}{}{}
und

\mavergleichskette
{\vergleichskette
{ Y
}
{ = }{ h^{-1}(c)
}
{  }{
}
{    }{
}
{   }{
}
}
{}{}{}
die
\definitionsverweis {Faser}{}{}
zu

\mavergleichskette
{\vergleichskette
{ c
}
{ \in }{ \R
}
{  }{
}
{    }{
}
{   }{
}
}
{}{}{,}
wobei $h$ in jedem Punkt von $Y$
\definitionsverweis {regulär}{}{}
sei. Ein auf einer offenen Umgebung

\mavergleichskette
{\vergleichskette
{ Y
}
{ \subseteq }{ U
}
{ \subseteq }{ \R^n
}
{    }{
}
{   }{
}
}
{}{}{}
definiertes
\definitionsverweis {stetiges}{}{}
\definitionsverweis {Vektorfeld}{}{}
\maabb {F} {U} {\R^n
} {}
mit

\mavergleichskettedisp
{\vergleichskette
{ F(P)
}
{ \in} { N_PY
}
{ } {
}
{ } {
}
{ } {
}
}
{}{}{}
für alle

\mavergleichskette
{\vergleichskette
{ P
}
{ \in }{ Y
}
{  }{
}
{    }{
}
{   }{
}
}
{}{}{}
heißt
\definitionswort {Normalenfeld}{}
auf $Y$.

}

Wenn dabei zusätzlich

\mavergleichskettedisp
{\vergleichskette
{ \Vert {N(P)} \Vert
}
{ =} { 1
}
{ } {
}
{ } {
}
{ } {
}
}
{}{}{}
gilt, so spricht man von einem \stichwort {Einheitsnormalenfeld} {.} Da die Normalengerade im betrachtenten Hyperflächenfall eindimensional ist, gibt es für jeden Punkt

\mavergleichskette
{\vergleichskette
{ P
}
{ \in }{ Y
}
{  }{
}
{    }{
}
{   }{
}
}
{}{}{}
nur zwei mögliche Werte für ein Einheitsnormalenfeld an diesem Punkt, die zueinander negativ sind. Dabei interessiert man sich hauptsächlich nur für die Werte des Feldes auf $Y$, man betrachtet also zwei Normalenfelder zu $Y$ als identisch, wenn sie auf $Y$ übereinstimmen. Die offene Umgebung ist nur nötig, um von differenzierbar sprechen zu können.

__NOEDITSECTION__

\inputfaktbeweis
{Hyperfläche/Glatt/Normalenfeld/Fakt}
{Lemma}
{}
{

\faktsituation {Es sei

\mavergleichskette
{\vergleichskette
{ W
}
{ \subseteq }{ \R^n
}
{  }{
}
{    }{
}
{   }{
}
}
{}{}{}
\definitionsverweis {offen}{}{,}
\maabb {h} { W } { \R
} {}
eine
\definitionsverweis {stetig differenzierbare Funktion}{}{}
und

\mavergleichskette
{\vergleichskette
{ Y
}
{ = }{ h^{-1}(c)
}
{  }{
}
{    }{
}
{   }{
}
}
{}{}{}
die
\definitionsverweis {Faser}{}{}
zu

\mavergleichskette
{\vergleichskette
{ c
}
{ \in }{ \R
}
{  }{
}
{    }{
}
{   }{
}
}
{}{}{,}
wobei $h$ in jedem Punkt von $Y$
\definitionsverweis {regulär}{}{}
sei.}
\faktfolgerung {Dann ist
\mathl{{ \frac{
\operatorname{Grad} \,  h  }{  \Vert {
\operatorname{Grad} \,  h } \Vert  } }}{} ein
\definitionsverweis {Einheitsnormalenfeld}{}{}
auf $Y$.}
\faktzusatz {}
\faktzusatz {}

}
{

Nach Voraussetzung ist das
\definitionsverweis {Gradientenfeld}{}{}
$\operatorname{Grad} \,  h$ nullstellenfrei auf $Y$ und daher wegen der vorausgesetzten Stetigkeit auch nullstellenfrei in einer offenen Umgebung

\mavergleichskette
{\vergleichskette
{ Y
}
{ \subseteq }{ U
}
{ \subseteq }{ \R^n
}
{    }{
}
{   }{
}
}
{}{}{.}
Das angegebene Vektorfeld ist also in einer offenen Umgebung von $Y$ definiert. Die Orthogonalität zu den Tangentialräumen an die Faser und die Normiertheit sind klar.

}

\inputbeispiel{}
{

\bild{ \begin{center}
\includegraphics[width=5.5cm]{ \bildeinlesungpng {Hyperboloid1} {png} }
\end{center}
\bildtext {Ein \stichwort {einschaliges Hyperboloid} {.}} }

\bildlizenz { Hyperboloid1.png } {} {RokerHRO} {Commons} {gemeinfrei} {}

Wir betrachten die Funktion
\maabbele {h} {\R^3} {\R
} {(x,y,z)} { x^2+y^2-z^2-1
} {}
und dazu die Faser über $0$, also das  \stichwort {einschalige Hyperboloid} {} $Y$. Der
\definitionsverweis {Gradient}{}{}
zu $h$ in einem Punkt
\mathl{(x,y,z)}{}  ist durch
\mathl{\begin{pmatrix}  2x \\ 2y\\  -2z   \end{pmatrix}}{}
und das
\definitionsverweis {totale Differential}{}{}
ist entsprechend durch
\mathl{\left(  2x  , \,   2y  , \,   -2z                 \right)}{} gegeben, daher ist $h$ in jedem Punkt von $Y$
\definitionsverweis {regulär}{}{.}
Der
\definitionsverweis {Tangentialraum}{}{}
in einem Punkt ist der Kern des totalen Differentials, eine Basis ist
\mathl{\begin{pmatrix}  y  \\ -x\\  0   \end{pmatrix} ,  \begin{pmatrix}  z  \\ 0 \\  x   \end{pmatrix}}{}
\zusatzklammer {bei

\mavergleichskettek
{\vergleichskettek
{ x
}
{ = }{ 0
}
{  }{
}
{    }{
}
{   }{
}
}
{}{}{}
muss man den zweiten Vektor durch $\begin{pmatrix}  0  \\ z \\  y   \end{pmatrix}$ ersetzen} {} {.}
Das
\definitionsverweis {Einheitsnormalenfeld}{}{}
ist

\mavergleichskettealign
{\vergleichskettealign
{ N( \begin{pmatrix}  x  \\ y \\ z   \end{pmatrix} )
}
{ =} { { \frac{
\operatorname{Grad} \, h(x,y,z)  }{  \Vert {
\operatorname{Grad} \, h(x,y,z) } \Vert  } }
}
{ =} { \left(  x^2+y^2+z^2  \right)^{ - { \frac{   1  }{  2 } } } \begin{pmatrix}  x  \\ y \\  -z   \end{pmatrix}
}
{ } {
}
{ } {
}
}
{}
{}{.}

}

\inputdefinition
{}
{

Es sei

\mavergleichskette
{\vergleichskette
{ W
}
{ \subseteq }{ \R^n
}
{  }{
}
{    }{
}
{   }{
}
}
{}{}{}
\definitionsverweis {offen}{}{,}
\maabb {h} { W } { \R
} {}
eine
\definitionsverweis {stetig differenzierbare Funktion}{}{}
und

\mavergleichskette
{\vergleichskette
{ Y
}
{ = }{ h^{-1}(c)
}
{  }{
}
{    }{
}
{   }{
}
}
{}{}{}
die
\definitionsverweis {Faser}{}{}
zu

\mavergleichskette
{\vergleichskette
{ c
}
{ \in }{ \R
}
{  }{
}
{    }{
}
{   }{
}
}
{}{}{,}
wobei $h$ in jedem Punkt von $Y$
\definitionsverweis {regulär}{}{}
sei. Dann nennt man die Fixierung eines
\definitionsverweis {Einheitsnormalenfeldes}{}{}
auf $Y$ eine
\definitionswort {Orientierung}{}
von $Y$.

}

Zu einer Orientierung gibt es stets die negative oder die entgegengesetzte Orientierung. Wenn eine beschreibende Funktion $h$ für $Y$ fixiert ist, so erhält man mit
[[Hyperfläche/Glatt/Normalenfeld/Fakt|Lemma 2.4]]
direkt die Orientierung
\mathl{{ \frac{
\operatorname{Grad} \,  h  }{  \Vert {
\operatorname{Grad} \,  h } \Vert  } }}{.} Allerdings gehört zu $-h$ die entgegengesetzte Orientierung, und es gibt keine kanonische Möglichkeit, eine davon auszuwählen. Bei einem nullstellenfreien Normalenfeld kann man stets zum zugehörigen Einheitsnormalenfeld übergehen und erhält somit eine Orientierung.

__NOEDITSECTION__

\inputfaktbeweis
{Hyperfläche/Glatt/Zusammenhängend/Orientierung/Fakt}
{Lemma}
{}
{

\faktsituation {Es sei

\mavergleichskette
{\vergleichskette
{ W
}
{ \subseteq }{ \R^n
}
{  }{
}
{    }{
}
{   }{
}
}
{}{}{}
\definitionsverweis {offen}{}{,}
\maabb {h} { W } { \R
} {}
eine
\definitionsverweis {stetig differenzierbare Funktion}{}{}
und

\mavergleichskette
{\vergleichskette
{ Y
}
{ = }{ h^{-1}(c)
}
{  }{
}
{    }{
}
{   }{
}
}
{}{}{}
die
\definitionsverweis {Faser}{}{}
zu

\mavergleichskette
{\vergleichskette
{ c
}
{ \in }{ \R
}
{  }{
}
{    }{
}
{   }{
}
}
{}{}{,}
wobei $h$ in jedem Punkt von $Y$
\definitionsverweis {regulär}{}{}
sei. Es sei $Y$
\definitionsverweis {zusammenhängend}{}{.}}
\faktfolgerung {Dann gibt es genau zwei
\definitionsverweis {Orientierungen}{}{}
auf $Y$.}
\faktzusatz {}
\faktzusatz {}

}
{

Nach
[[Hyperfläche/Glatt/Normalenfeld/Fakt|Lemma 2.4]]
gibt es ein Einheitsnormalenfeld, das eine Orientierung repräsentiert, und die Negation davon liefert eine weitere Orientierung. Diese nennen wir $G$ bzw. $-G$. Es sei nun
\maabbdisp {N} { Y } { \R^n
} {}
ein beliebiges stetiges Einheitsnormalenfeld. Wir betrachten die Übereinstimmungsorte

\mavergleichskettedisp
{\vergleichskette
{ Y_1
}
{ =} { \left\{ P \in Y  \mid   N(P) = G(P)         \right\}
}
{ } {
}
{ } {
}
{ } {
}
}
{}{}{}
und

\mavergleichskettedisp
{\vergleichskette
{ Y_2
}
{ =} { \left\{ P \in Y  \mid   N(P) = - G(P)         \right\}
}
{ } {
}
{ } {
}
{ } {
}
}
{}{}{.}
Da es für jeden Punkt

\mavergleichskette
{\vergleichskette
{ P
}
{ \in }{ Y
}
{  }{
}
{    }{
}
{   }{
}
}
{}{}{}
nur die beiden möglichen Werte

\mavergleichskettedisp
{\vergleichskette
{ N(P)
}
{ =} { \pm G(P)
}
{ } {
}
{ } {
}
{ } {
}
}
{}{}{}
gibt, ist $Y$ die disjunkte Vereinigung von
\mathkor {} {Y_1} {und} {Y_2} {.}
Als Übereinstimmungsort von stetigen Funktionen sind
\mathkor {} {Y_1} {und} {Y_2} {}
\definitionsverweis {abgeschlossen}{}{}
\zusatzklammer {und damit auch
\definitionsverweis {offen}{}{}
in $Y$} {} {.}
Aufgrund des Zusammenhangs ist also

\mavergleichskette
{\vergleichskette
{ Y
}
{ = }{ Y_1
}
{  }{
}
{    }{
}
{   }{
}
}
{}{}{}
oder

\mavergleichskette
{\vergleichskette
{ Y
}
{ = }{ Y_2
}
{  }{
}
{    }{
}
{   }{
}
}
{}{}{.}

}

  \zwischenueberschrift{Die Gauß-Abbildung}

\inputdefinition
{}
{

Es sei

\mavergleichskette
{\vergleichskette
{ W
}
{ \subseteq }{ \R^n
}
{  }{
}
{    }{
}
{   }{
}
}
{}{}{}
\definitionsverweis {offen}{}{,}
\maabb {h} { W } { \R
} {}
eine
\definitionsverweis {stetig differenzierbare Funktion}{}{}
und

\mavergleichskette
{\vergleichskette
{ Y
}
{ = }{ h^{-1}(c)
}
{  }{
}
{    }{
}
{   }{
}
}
{}{}{}
die
\definitionsverweis {Faser}{}{}
zu

\mavergleichskette
{\vergleichskette
{ c
}
{ \in }{ \R
}
{  }{
}
{    }{
}
{   }{
}
}
{}{}{,}
wobei $h$ in jedem Punkt von $Y$
\definitionsverweis {regulär}{}{}
sei. Es sei eine
\definitionsverweis {Orientierung}{}{}
auf $Y$ fixiert. Dann heißt die Abbildung
\maabbdisp {} { Y } { S^{n-1}
} {,}
die jeden Punkt

\mavergleichskette
{\vergleichskette
{ P
}
{ \in }{ Y
}
{  }{
}
{    }{
}
{   }{
}
}
{}{}{}
auf seinen durch die Orientierung fixierten Einheitsnormalenvektor abbildet, die
\definitionswort {Gauß-Abbildung}{}
zu $Y$.

}

Im Allgemeinen denken wir uns Tangentialvektoren und Normalenvektoren an den Punkt

\mavergleichskette
{\vergleichskette
{ P
}
{ \in }{ Y
}
{  }{
}
{    }{
}
{   }{
}
}
{}{}{}
angeheftet, hier ist es aber wichtig, den Normalenvektor als einen Punkt auf der \anfuehrung{neutralen}{} $n-1$- dimensionalen Sphäre zu betrachten. Die Gauß-Abbildung eröffnet eine Möglichkeit, eine beliebige glatte Hyperfläche mit der besonders einfachen Hyperfläche, nämlich der Kugeloberfläche, in Beziehung zu setzen.

__NOEDITSECTION__

\inputfaktbeweis
{Hyperfläche/Glatt/Gauß-Abbildung/Eigenschaften/Fakt}
{Lemma}
{}
{

\faktsituation {Es sei

\mavergleichskette
{\vergleichskette
{ W
}
{ \subseteq }{ \R^n
}
{  }{
}
{    }{
}
{   }{
}
}
{}{}{}
\definitionsverweis {offen}{}{,}
\maabb {h} { W } { \R
} {}
eine
\definitionsverweis {stetig differenzierbare Funktion}{}{}
und

\mavergleichskette
{\vergleichskette
{ Y
}
{ = }{ h^{-1}(c)
}
{  }{
}
{    }{
}
{   }{
}
}
{}{}{}
die
\definitionsverweis {Faser}{}{}
zu

\mavergleichskette
{\vergleichskette
{ c
}
{ \in }{ \R
}
{  }{
}
{    }{
}
{   }{
}
}
{}{}{,}
wobei $h$ in jedem Punkt von $Y$
\definitionsverweis {regulär}{}{}
sei. Es sei eine
\definitionsverweis {Orientierung}{}{}
auf $Y$ fixiert.}
\faktuebergang {Dann gelten für die
\definitionsverweis {Gauß-Abbildung}{}{}
folgende Aussagen.}
\faktfolgerung {\aufzaehlungvier{Die Gauß-Abbildung ist
\definitionsverweis {stetig}{}{.}
}{Bei

\mavergleichskette
{\vergleichskette
{ Y
}
{ = }{ S^{n-1}
}
{  }{
}
{    }{
}
{   }{
}
}
{}{}{}
ist die Gauß-Abbildung
\zusatzklammer {wenn die Orientierung nach außen zeigt} {} {}
die Identität oder die antipodale Abbildung.
}{Wenn $Y$
\definitionsverweis {zusammenhängend}{}{}
ist, so sind die Gauß-Abbildungen zu den beiden Orientierungen antipodal zueinander.
}{Sei $Y$ zusammenhängend. Dann ist die Gauß-Abbildung genau dann konstant, wenn $Y$ ein offener Ausschnitt aus einem
\definitionsverweis {affin-linearen Unterraum}{}{}
der Dimension $n-1$ ist.
}}
\faktzusatz {}
\faktzusatz {}

}
{

\aufzaehlungvier{Dies folgt aus
[[Hyperfläche/Glatt/Normalenfeld/Fakt|Lemma 2.4]].
}{Klar.
}{Klar.
}{Die Rückrichtung ist klar. Zum Beweis der Hinrichtung sei der Einheitsnormalenvektor konstant gleich

\mavergleichskettedisp
{\vergleichskette
{ v
}
{ =} { \begin{pmatrix} a_1  \\\vdots\\ a_n   \end{pmatrix}
}
{ } {
}
{ } {
}
{ } {
}
}
{}{}{.}
Dann ist

\mavergleichskettedisp
{\vergleichskette
{
\operatorname{Grad} \,  h  (  P )
}
{ =} { g(P) v
}
{ } {
}
{ } {
}
{ } {
}
}
{}{}{}
mit einer stetigen nullstellenfreien Funktion
\maabb {g} { Y } {\R
} {.}
Wir können $h$ durch
\mathl{{ \frac{   h  }{  \Vert { h } \Vert  } }}{} ersetzen ohne $Y$ zu verändern. Dann ist nach
[[Hyperfläche/Glatt/Zusammenhängend/Orientierung/Fakt|Lemma 2.7]]

\mavergleichskettedisp
{\vergleichskette
{
\operatorname{Grad} \,  h  (  P )
}
{ =} { \pm v
}
{ } {
}
{ } {
}
{ } {
}
}
{}{}{}
und daher

\mavergleichskettedisp
{\vergleichskette
{ h(x_1 , \ldots , x_n)
}
{ =} { \pm \sum_{ i = 1}^n a_ix_i
}
{ } {
}
{ } {
}
{ } {
}
}
{}{}{.}
Die Faser zu $h$ auf ganz $\R^n$ ist ein affin-linearer Untervektorraum.}

}

Wir erinnern an den folgenden Satz.

\inputfakt{Hyperfläche/Glatt und kompakt/Jede Hyperebene als Tangentialraum/Fakt}{Satz}{}
{

\faktsituation {Es sei

\mavergleichskette
{\vergleichskette
{ U
}
{ \subseteq }{\R^n
}
{  }{
}
{    }{
}
{   }{
}
}
{}{}{}
eine
\definitionsverweis {offene Teilmenge}{}{}
und sei
\maabbdisp {f} { U } {\R
} {}
eine
\definitionsverweis {stetig differenzierbare Funktion}{}{.}}
\faktvoraussetzung {Die
\definitionsverweis {Faser}{}{}

\mavergleichskette
{\vergleichskette
{M
}
{ = }{ f^{-1}(b)
}
{  }{
}
{    }{
}
{   }{
}
}
{}{}{}
von $f$ zu einem Punkt

\mavergleichskette
{\vergleichskette
{b
}
{ \in }{\R
}
{  }{
}
{    }{
}
{   }{
}
}
{}{}{}
sei
\definitionsverweis {kompakt}{}{}
und in jedem Punkt
\definitionsverweis {regulär}{}{.}}
\faktfolgerung {Dann ist jeder
\mathl{(n-1)}{-}dimensionale Unterraum

\mavergleichskette
{\vergleichskette
{V
}
{ \subset }{ \R^n
}
{  }{
}
{    }{
}
{   }{
}
}
{}{}{}
für mindestens einen Punkt

\mavergleichskette
{\vergleichskette
{a
}
{ \in }{ M
}
{  }{
}
{    }{
}
{   }{
}
}
{}{}{}
gleich dem
\definitionsverweis {Tangentialraum}{}{}

\mathl{T_aM}{.}}
\faktzusatz {}
\faktzusatz {}

}
__NOEDITSECTION__

\inputfaktbeweis
{Hyperfläche/Glatt und kompakt/Gauß-Abbildung surjektiv/Fakt}
{Satz}
{}
{

\faktsituation {Es sei

\mavergleichskette
{\vergleichskette
{ W
}
{ \subseteq }{\R^n
}
{  }{
}
{    }{
}
{   }{
}
}
{}{}{}
eine
\definitionsverweis {offene Teilmenge}{}{}
und sei
\maabbdisp {h} { W } {\R
} {}
eine
\definitionsverweis {stetig differenzierbare Funktion}{}{.}}
\faktvoraussetzung {Die
\definitionsverweis {Faser}{}{}

\mavergleichskette
{\vergleichskette
{ Y
}
{ = }{ h^{-1}(c)
}
{  }{
}
{    }{
}
{   }{
}
}
{}{}{}
von $h$ zu

\mavergleichskette
{\vergleichskette
{ c
}
{ \in }{\R
}
{  }{
}
{    }{
}
{   }{
}
}
{}{}{}
sei
\definitionsverweis {kompakt}{}{}
und in jedem Punkt
\definitionsverweis {regulär}{}{.}}
\faktfolgerung {Dann ist die
\definitionsverweis {Gauß-Abbildung}{}{}
\maabbdisp {} { Y } { S^{n-1}
} {}
surjektiv.}
\faktzusatz {}
\faktzusatz {}

}
{

Ein Einheitsnormalenvektor beschreibt über die Orthogonalitätsrelation eine Hyperebene, also einen $n-1$-dimensionalen Untervektorraum, wobei auch der negierte Einheitsnormalenvektor die gleiche Hyperebene beschreibt. [[Hyperfläche/Glatt und kompakt/Jede Hyperebene als Tangentialraum/Fakt|Satz 2.10]]
zeigt, dass jede Hyperebene als ein Tangentialraum von $Y$ auftritt. Der Beweis von
[[Hyperfläche/Glatt und kompakt/Jede Hyperebene als Tangentialraum/Fakt|Satz 2.10]]
zeigt aber ferner
\zusatzklammer {wenn man dort neben dem Maximum auch das Minimum betrachtet} {} {,}
dass beide Normaleneinheitsvektoren auftreten.

}

 [[Kategorie:Latexseite]]
